%%==============================================================================
%% 第一章 绪论（示例占位内容，请替换）
%%==============================================================================
\chapter{绪论}\label{chap:intro}

\section{研究背景与意义}
本节为示例占位内容，用于演示正文排版。正文采用宋体小四号，行距为固定值约二十磅，
段落首行缩进两个字符。请将此处替换为研究背景与意义的实际论述。

正文中引用参考文献时，使用上标数字标注，例如已有研究\cite{ref-journal}对相关问题进行了
讨论，另有若干工作\cite{ref-book,ref-thesis}也有所涉及。请将本段替换为实际内容。

\section{国内外研究现状}

\subsection{某方面研究现状}
本小节为示例占位内容。请将此处替换为某一方面研究现状的综述，并在适当位置标注
参考文献\cite{ref-conf}。

\subsection{某方面研究现状}
本小节为示例占位内容。请将此处替换为另一方面研究现状的综述\cite{ref-en}。

\section{本文主要研究内容}
本文的章节安排如下：
\begin{enumerate}[label={（\arabic*）},leftmargin=*]
  \item 第二章介绍相关理论与研究方法；
  \item 第三章给出实验设计与结果分析；
  \item 第四章总结全文并展望后续工作。
\end{enumerate}
